% =============================================================================
% New paper (separate from paper/): skeleton only.
% Single .tex file. Figures go in figures/. Bibliography in refs.bib.
%
% Method assumed: CA seed with uniform module weights -> low-resolution
% consensus of negative-contrast gates (internally MK4-LR) -> fine-only 64x64
% self-attention propagation (internally V0) -> threshold (+ DenseCRF).
% No capped hierarchy, no SA-guided module weights.
% The freeze of this pipeline is still an owner decision (see runs/42).
%
% Conventions
%   \x          = empty table cell to fill.
%   \citeneeded = citation still to add. Take bib entries from the PDFs in
%                 "relevant literature/" and check every entry against them;
%                 the entries copied from paper/references.bib are unverified.
%   % SRC:      = where the numbers should come from (run, or NEW RUN).
%   Manuscript text must not use internal names (A3, V3, V0, MK4, LR, SGxx).
%
% Format: Springer LNCS, used by recent ICPR editions. Check the ICPR 2027
% call for the class, page limit and whether references count.
% =============================================================================
\documentclass[runningheads]{llncs}

\usepackage[T1]{fontenc}
\usepackage{graphicx}
\usepackage{booktabs}
\IfFileExists{multirow.sty}{\usepackage{multirow}}{\newcommand{\multirow}[3]{##3}} % local TeX lacks multirow
\usepackage{amsmath,amssymb}
\usepackage{xcolor}
\usepackage[hidelinks]{hyperref}
\graphicspath{{figures/}}

\newcommand{\method}{\textsc{Method}}          % name to be decided
\newcommand{\x}{\textcolor{gray}{--}}          % empty cell
\newcommand{\citeneeded}{\textcolor{red}{[?]}} % missing citation
\newcommand{\figbox}[2]{\fbox{\parbox[c][#2][c]{0.95\linewidth}{\centering\small\textcolor{gray}{#1}}}}

\begin{document}

\title{Title}
\titlerunning{Title}
\author{Author(s)}
\authorrunning{Author(s)}
\institute{Institution}

\maketitle

\begin{abstract}
\end{abstract}

\keywords{Keyword \and Keyword \and Keyword}

% =============================================================================
\section{Introduction}
\label{sec:intro}
% Problem, gap, idea, contributions, headline numbers.

\begin{figure}[t]
\centering
\figbox{Teaser: photo, class word, cross-attention seed, gate, final mask; one seen and one unseen class.}{3.2cm}
% FILE: figures/teaser.pdf
\caption{}
\label{fig:teaser}
\end{figure}

% =============================================================================
\section{Related Work}
\label{sec:related}

\subsection{Segmentation from diffusion attention}
% DAAM, OVAM, DiffSegmenter, iSeg, SeeDiff, M2N2, FA-Seg, Meng et al. 2026.
% FA-Seg (Fig. 3) already shows SD attends to absent classes; their fix is
% distractor classes in the prompt.

\subsection{Prototype and reference-set methods}
% Position the novelty here.
% OVDiff: per-class positive and negative (background-context) prototypes from
%   generated support images, k-means, nearest prototype. Closest for negatives;
%   ours = the word's own false activations on photos without the word, used as
%   a gate on the attention map, not as the classifier.
% MaViSeg: prototypes from the top 2% tokens of unlabelled real images, in the
%   model's own features, positives only, DiT backbones. Closest for profiles.
% ReCo: retrieved real images; hand-picked context classes as a 1-P gate.
% FreeDA, FOSSIL: caption-driven generated references with DINO features.
% Still to obtain: ConceptAttention, Seg4Diff, CLIP as RNN (CaR).

\subsection{Training-free open-vocabulary segmentation with CLIP}
% MaskCLIP, SCLIP, ClearCLIP, NACLIP, ProxyCLIP, PEARL.

\subsection{Propagation and refinement}
% Random walks (NERVE, SPARK), self-attention propagation, DenseCRF.

% =============================================================================
\section{Method}
\label{sec:method}

\begin{figure}[t]
\centering
\figbox{Method overview: one noising pass; cross-attention seed; reference profiles from captioned photos; per-module negative-contrast gates; low-resolution consensus; fine self-attention propagation; threshold.}{4cm}
% FILE: figures/method.pdf
\caption{}
\label{fig:method}
\end{figure}

\subsection{Attention capture}
\label{sec:capture}
% One SD1.5 pass at t = 100, prompt template, image at 512.

\subsection{Cross-attention seed}
\label{sec:seed}
% Uniform module weights.

\subsection{Reference profiles from captioned photos}
\label{sec:profiles}
% Caption matching, object profile (top 5% queries), negatives (top 1% on
% non-mentioning photos, 4 spherical k-means centroids). No labels.

\subsection{Negative-contrast gate}
\label{sec:gate}
% Nearest-negative rule per module, gamma.

\subsection{Low-resolution consensus}
\label{sec:consensus}
% Mean of the six gates below 64x64, upsampled to 64x64, applied at every module.

\subsection{Propagation and read-out}
\label{sec:propagation}
% Two 64x64 self-attention steps, threshold, optional DenseCRF.

% =============================================================================
\section{Experimental Protocol}
\label{sec:protocol}

\subsection{Data}
\label{sec:data}

\begin{table}[t]
\centering
\caption{Datasets and their role. Items are (image, class) pairs scored as binary masks; on the real evaluation sets, only items whose class covers at least 1\% of the image are counted (primary items). Profile photos are selected by their captions only; no masks or labels are used.}
% Evaluation sets and their role. One row per set used anywhere in the paper.
% SRC: runs/43 (generated), runs/32-37 (VOC), runs/40 items_train (threshold check: 1,200 items, 993 photos),
%      runs/42 items_val (val2017 minus Section 136 set, primary items) and items_prof (profile pool, 100 + 480 per word).
% Roles of VOC train / COCO train threshold set depend on the Table 3 run plan; update if that plan changes them.
\label{tab:sets}
\footnotesize\setlength{\tabcolsep}{4pt}
\begin{tabular}{@{}llcccl@{}}
\toprule
Set & Images & Classes & Photos & Items & Role \\
\midrule
VOC-Sim~\cite{marcosmanchon2024ovam}  & generated & 20 & 490 & 490 & evaluation (Tab.~\ref{tab:generated}) \\
COCO-Cap~\cite{marcosmanchon2024ovam} & generated & 20 & 707 & 707 & evaluation (Tab.~\ref{tab:generated}) \\
PASCAL VOC 2012 train\citeneeded      & real      & 20 & 1,464 & 2,170 & development only \\
PASCAL VOC 2012 val\citeneeded        & real      & 20 & 1,409 & 1,947 & evaluation (Tab.~\ref{tab:real}) \\
COCO train2017 (threshold check)\citeneeded & real & 20 & 993 & 1,200 & threshold check \\
COCO train2017 (profiles, captions only) & real   & 80 & 7,986 & -- & reference profiles \\
PASCAL Context val (59 classes)\citeneeded & real & 59 & 5,104 & \x & evaluation (Tab.~\ref{tab:real}) \\
COCO-Object (COCO val2017)\citeneeded & real & 80 & 4,571 & 9,816 & evaluation (Tab.~\ref{tab:real}) \\
\quad 20 VOC classes                  & real      & 20 & 3,582 & 5,414 & seen in development (Tab.~\ref{tab:openvocab}) \\
\quad other 60 classes                & real      & 60 & 3,038 & 4,402 & never used in development (Tab.~\ref{tab:openvocab}) \\
\bottomrule
\end{tabular}
\end{table}

\subsection{Metrics and statistics}
\label{sec:metrics}
% Class-macro mIoU on (image, class) items; >=1% area primary;
% paired photo bootstrap, 2,000 resamples, 95% intervals.

\subsection{Selection rules and pre-registration}
\label{sec:prereg}
% What was fixed before scoring; development vs confirmation sets.

\subsection{Baselines}
\label{sec:baselines}
% Which methods we re-run under the same protocol, and which numbers are quoted.

\subsection{Implementation details}
\label{sec:impl}

% =============================================================================
\section{Results}
\label{sec:results}

\subsection{Generated benchmarks}
\label{sec:res-generated}

\begin{table}[t]
\centering
\caption{Class-word segmentation on the generated benchmarks (mIoU, \%). The last two columns split COCO-Cap by whether the queried class word appears in the generation prompt (465 / 242 images). Published rows are quoted from~\cite{marcosmanchon2024ovam} and report only the full sets; $^\dagger$~re-run by us, as released, on the same images as our method. Our maps use a fixed threshold of 0.55, determined empirically in earlier experiments and not tuned here; bold marks the best method that uses no annotated masks.}
% Comparison with published class-word results on the OVAM benchmarks (mIoU).
% Last two columns: COCO-Cap split by whether the class word is in the
% generation prompt (OVAM's word_included; the OVAM sense of querying a word
% not in the prompt). Published rows report only the full sets.
% $^\dagger$ = re-run by us on the same regenerated images, as released
%   (NERVE: SD 2.1-base, own binarisation; SCLIP: ViT-B/16 VOC21 prediction).
% Ours: fixed threshold 0.55 ("determined empirically in earlier experiments",
%   as OVAM words it; it came from VOC-Sim for the earlier pipeline: say so).
%   Profiles from the other half of each set (no image in its own profile).
% SRC: published rows quoted from OVAM Table 1; all other rows runs/43.
\label{tab:generated}
\footnotesize\setlength{\tabcolsep}{3pt}
\begin{tabular}{@{}lccccc@{}}
\toprule
 & & & & \multicolumn{2}{c}{word in prompt} \\
\cmidrule(lr){5-6}
Method & Training & VOC-Sim & COCO-Cap & in & not in \\
\midrule
DAAM~\cite{tang2023daam}                 & no  & 66.2 & 48.4 & -- & -- \\
Attn2Mask~\cite{yoshihashi2024attn2mask} & no  & 68.7 & 55.0 & -- & -- \\
OVAM~\cite{marcosmanchon2024ovam}        & no  & 70.4 & 58.2 & -- & -- \\
NERVE~\cite{mahatha2025nerve}$^\dagger$  & no  & \textbf{80.5} & 70.2 & 73.8 & 61.5 \\
SCLIP~\cite{wang2024sclip}$^\dagger$     & no  & 71.9 & 58.3 & 61.6 & 51.4 \\
\midrule
Grounded Diffusion~\cite{li2023groundeddiffusion} & yes & 62.1 & 50.2 & -- & -- \\
DatasetDM~\cite{wu2023datasetdm}         & yes & 80.3 & 59.3 & -- & -- \\
OVAM + token opt.~\cite{marcosmanchon2024ovam} & 1 mask/cls & 82.5 & 69.2 & -- & -- \\
\midrule
No gate (ours, ablation)                 & no  & 78.2 & 66.2 & 67.3 & 63.9 \\
\method{} (raw)                          & no  & 80.1 & 73.6 & 74.6 & 71.7 \\
\method{} + DenseCRF                     & no  & \textbf{80.5} & \textbf{74.4} & \textbf{75.3} & \textbf{72.8} \\
\bottomrule
\end{tabular}
\end{table}

\subsection{Real photographs}
\label{sec:res-real}

\begin{table}[t]
\centering
\caption{}
% Real photographs, class word given, binary mask per (image, class) item.
% Columns: the three object benchmarks reported by nearly all diffusion-based methods (OVDiff, iSeg,
%   DiffSegmenter, FreeDA; also the first three of NERVE/NACLIP/SCLIP). Their datasets, OUR protocol
%   (class word given), so published multi-class numbers are NOT comparable; every row is re-run (runs/44).
% Every row re-run by us with the same items, thresholds rule and bootstrap.
% SRC: NEW RUN for all baselines on real photos; ours from the frozen capture
%      (runs/40 and runs/42 use the relevant pipelines but different seeds/arms).
\label{tab:real}
\footnotesize\setlength{\tabcolsep}{3.5pt}
\begin{tabular}{@{}llccc@{}}
\toprule
 & & PASCAL VOC & PASCAL Context & COCO-Object \\
Family & Method & (20) & (59) & (80) \\
\midrule
\multirow{4}{*}{Diffusion attention}
 & DAAM~\cite{tang2023daam}            & \x & \x & \x \\
 & OVAM (no tok.\ opt.)~\cite{marcosmanchon2024ovam} & \x & \x & \x \\
 & DiffSegmenter\citeneeded            & \x & \x & \x \\
 & iSeg\citeneeded                     & \x & \x & \x \\
\midrule
\multirow{4}{*}{CLIP}
 & SCLIP~\cite{wang2024sclip}                    & \x & \x & \x \\
 & ClearCLIP\citeneeded                & \x & \x & \x \\
 & NACLIP\citeneeded                   & \x & \x & \x \\
 & ProxyCLIP\citeneeded                & \x & \x & \x \\
\midrule
Random walk & NERVE~\cite{mahatha2025nerve} & \x & \x & \x \\
\midrule
Prototypes & OVDiff~\cite{karazija2024ovdiff} & \x & \x & \x \\
\midrule
\multirow{2}{*}{Ours}
 & \method{} (raw)                     & \x & \x & \x \\
 & \method{} + DenseCRF                & \x & \x & \x \\
\bottomrule
\end{tabular}
\end{table}

\subsection{Standard multi-class benchmarks}
\label{sec:res-multiclass}

\begin{table}[t]
\centering
\caption{}
% Standard training-free open-vocabulary semantic segmentation (mIoU), all
% classes compete per pixel. Published numbers quoted as reported.
% OPTIONAL (owner decision 2026-11-24): the paper uses OVAM's sense of open
% vocabulary and does not claim open-vocabulary segmentation. Keep only if a
% multi-class rule is added later; otherwise move to supplementary or drop.
% SRC: NEW RUN; needs the multi-class competition rule (not designed).
% $^\ddagger$ = diffusion-transformer backbone (FLUX), not SD1.5; quote only.
\label{tab:multiclass}
\footnotesize\setlength{\tabcolsep}{4pt}
\begin{tabular}{@{}lccc@{}}
\toprule
Method & VOC20 & VOC21 & COCO-Object \\
\midrule
MaskCLIP\citeneeded                 & \x & \x & \x \\
SCLIP~\cite{wang2024sclip}                    & \x & \x & \x \\
ClearCLIP\citeneeded                & \x & \x & \x \\
NACLIP\citeneeded                   & \x & \x & \x \\
ProxyCLIP\citeneeded                & \x & \x & \x \\
PEARL~\cite{pei2026pearl}           & \x & \x & \x \\
DiffSegmenter\citeneeded            & \x & \x & \x \\
iSeg\citeneeded                     & \x & \x & \x \\
Meng et al.~\cite{meng2026making}   & \x & \x & \x \\
OVDiff~\cite{karazija2024ovdiff}    & \x & \x & \x \\
MaViSeg~\cite{chakraborty2026maviseg}$^\ddagger$ & \x & \x & \x \\
\midrule
\method{}                           & \x & \x & \x \\
\bottomrule
\end{tabular}
\end{table}

\begin{figure}[t]
\centering
\figbox{Qualitative comparison: rows = photos (seen and unseen classes, VOC and COCO); columns = photo, ground truth, two strongest baselines, ours.}{4.5cm}
% FILE: figures/qualitative.pdf
\caption{}
\label{fig:qualitative}
\end{figure}

\subsection{Classes never used in development}
\label{sec:res-openvocab}

\begin{table}[t]
\centering
\caption{}
% Gate effect on the 20 development classes and the 60 classes never used in
% development, by number of caption-matched reference photos per class.
% mIoU (raw, threshold from seen-class COCO train items); differences with
% 95% paired bootstrap intervals.
% SRC: runs/42 (U1, U2, S1, S2; size 20 primary; sizes 5 and 100 descriptive).
\label{tab:openvocab}
\footnotesize\setlength{\tabcolsep}{3.5pt}
\begin{tabular}{@{}lcccc@{}}
\toprule
 & \multicolumn{2}{c}{VOC classes (20)} & \multicolumn{2}{c}{Other COCO classes (60)} \\
\cmidrule(lr){2-3}\cmidrule(lr){4-5}
Gate & mIoU & $\Delta$ vs none & mIoU & $\Delta$ vs none \\
\midrule
None                                & \x & --  & \x & -- \\
Per-module gate, 20 photos          & \x & \x  & \x & \x \\
Low-res.\ consensus, 5 photos       & \x & \x  & \x & \x \\
Low-res.\ consensus, 20 photos      & \x & \x  & \x & \x \\
Low-res.\ consensus, 100 photos     & \x & \x  & \x & \x \\
\bottomrule
\end{tabular}
\end{table}

\begin{figure}[t]
\centering
\figbox{Per-class gain over no gate for the 60 other COCO classes, against the number of reference photos; seen classes as a reference.}{4cm}
% FILE: figures/openvocab_classes.pdf
\caption{}
\label{fig:openvocab}
\end{figure}

% =============================================================================
\section{Analysis}
\label{sec:analysis}

\subsection{Components}
\label{sec:ablation}

\begin{table}[t]
\centering
\caption{}
% Component ablation, one capture per dataset, raw mIoU unless stated.
% The last block adds the two components of the earlier pipeline that we
% remove (show they do not help once the consensus gate is on).
% SRC: real sets from a NEW frozen capture; runs/37, /38, /40 give
%      exploratory/confirmatory precedents (with SA weights on).
\label{tab:ablation}
\footnotesize\setlength{\tabcolsep}{3pt}
\begin{tabular}{@{}lccccc@{}}
\toprule
 & & & \multicolumn{2}{c}{COCO val2017} & \\
\cmidrule(lr){4-5}
Configuration & VOC-Sim & COCO-Cap & VOC cl. & other & VOC val \\
\midrule
Cross-attention seed                      & \x & \x & \x & \x & \x \\
\quad + fine propagation                  & \x & \x & \x & \x & \x \\
\quad + per-module gate                   & \x & \x & \x & \x & \x \\
\quad + low-res.\ consensus (\method{}) & \x & \x & \x & \x & \x \\
\quad + DenseCRF                          & \x & \x & \x & \x & \x \\
\midrule
\multicolumn{6}{@{}l}{\emph{Added back to \method{} (removed from the earlier pipeline):}} \\
\quad + attention-guided weights & \x & \x & \x & \x & \x \\
\quad + coarse-to-fine hierarchy     & \x & \x & \x & \x & \x \\
\bottomrule
\end{tabular}
\end{table}

\subsection{Gate design}
\label{sec:gate-design}

\begin{table}[t]
\centering
\caption{}
% Gate design choices; everything else fixed. Raw mIoU.
% SRC: consensus variants: runs/39 (stage 1), runs/41 (end to end, median and
%      trimmed fail); negative type and K sweep: NEW RUN on the frozen pipeline,
%      all in one capture. The negative-type block carries the novelty claim:
%      none = positives only (as MaViSeg's prototypes);
%      context = background of the positive reference photos (as OVDiff);
%      false activations = the word's top attention on photos without it (ours).
\label{tab:gate}
\footnotesize\setlength{\tabcolsep}{4pt}
\begin{tabular}{@{}llccc@{}}
\toprule
 & & & \multicolumn{2}{c}{COCO val2017} \\
\cmidrule(lr){4-5}
Factor & Setting & VOC val & VOC cl. & other 60 \\
\midrule
\multirow{3}{*}{Negative type}
 & none (object profile only) & \x & \x & \x \\
 & context of reference photos & \x & \x & \x \\
 & false activations (ours)   & \x & \x & \x \\
\midrule
\multirow{3}{*}{Centroids}
 & 1                          & \x & \x & \x \\
 & 2                          & \x & \x & \x \\
 & 4 (ours)                   & \x & \x & \x \\
\midrule
\multirow{4}{*}{Modules}
 & each module separately     & \x & \x & \x \\
 & consensus, all modules     & \x & \x & \x \\
 & consensus, $64{\times}64$ only & \x & \x & \x \\
 & consensus, below $64{\times}64$ & \x & \x & \x \\
\midrule
\multirow{3}{*}{Combination}
 & mean                       & \x & \x & \x \\
 & median                     & \x & \x & \x \\
 & trimmed mean               & \x & \x & \x \\
\bottomrule
\end{tabular}
\end{table}

\subsection{Why low resolution}
\label{sec:why-lowres}

\begin{figure}[t]
\centering
\figbox{Object-vs-background separability (AUROC) of each module's gate, with 64$\times$64 modules pooled to low resolution; example gates per module.}{4cm}
% FILE: figures/gate_resolution.pdf
% SRC: runs/39 (exploratory).
\caption{}
\label{fig:why-lowres}
\end{figure}

% =============================================================================
\section{Limitations}
\label{sec:limitations}

\section{Conclusion}
\label{sec:conclusion}

% =============================================================================
\bibliographystyle{splncs04}
\bibliography{refs}

% =============================================================================
% SUPPLEMENTARY MATERIAL
% Kept in this file. Split into its own PDF at submission if ICPR requires it.
% =============================================================================
\clearpage
\appendix
\section*{Supplementary Material}

\section{Pre-registration record}
\label{sec:s-prereg}

\begin{table}[h]
\centering
\caption{}
% Every claim in the paper, the test behind it, when it was fixed, and status.
% SRC: runs/*/PREREG.md and result.md, commit hashes.
\label{tab:s-prereg}
\footnotesize\setlength{\tabcolsep}{3pt}
\begin{tabular}{@{}p{3.2cm}p{2.6cm}p{2cm}p{1.8cm}p{1.6cm}@{}}
\toprule
Claim & Test & Data & Status & Result \\
\midrule
\x & \x & \x & \x & \x \\
\x & \x & \x & \x & \x \\
\x & \x & \x & \x & \x \\
\bottomrule
\end{tabular}
\end{table}

\section{Hyperparameters}
\label{sec:s-hparams}

\begin{table}[h]
\centering
\caption{}
% Every setting of the frozen pipeline and where it was chosen.
\label{tab:s-hparams}
\footnotesize
\begin{tabular}{@{}lll@{}}
\toprule
Setting & Value & Chosen on \\
\midrule
Model, timestep, resolution        & \x & \x \\
Prompt template                    & \x & \x \\
Object-profile query fraction      & \x & \x \\
Negative query fraction, centroids & \x & \x \\
Gate exponent $\gamma$             & \x & \x \\
Consensus modules                  & \x & \x \\
Propagation steps                  & \x & \x \\
Threshold                          & \x & \x \\
DenseCRF parameters                & \x & \x \\
\bottomrule
\end{tabular}
\end{table}

\section{Reference profiles}
\label{sec:s-profiles}
% Caption-matching rules, ambiguous words, photos per class.

\begin{table}[h]
\centering
\caption{}
% Caption rule and number of matched reference photos per class (80 classes).
% SRC: runs/42 items/fit logs.
\label{tab:s-profiles}
\footnotesize
\begin{tabular}{@{}llc@{}}
\toprule
Class & Caption rule & Photos \\
\midrule
\x & \x & \x \\
\bottomrule
\end{tabular}
\end{table}

\section{Per-class results}
\label{sec:s-perclass}

\begin{table}[h]
\centering
\caption{}
% Per-class IoU, PASCAL VOC 2012 val; ours vs strongest baselines.
\label{tab:s-perclass-voc}
\footnotesize\setlength{\tabcolsep}{2pt}
\begin{tabular}{@{}lcccc@{}}
\toprule
Class & No gate & Per-module gate & \method{} & Best baseline \\
\midrule
\x & \x & \x & \x & \x \\
\bottomrule
\end{tabular}
\end{table}

\begin{table}[h]
\centering
\caption{}
% Per-class IoU, COCO val2017, all 80 classes, with reference-photo counts.
\label{tab:s-perclass-coco}
\footnotesize\setlength{\tabcolsep}{2pt}
\begin{tabular}{@{}lccccc@{}}
\toprule
Class & Photos & No gate & Per-module gate & \method{} & Best baseline \\
\midrule
\x & \x & \x & \x & \x & \x \\
\bottomrule
\end{tabular}
\end{table}

\section{Sensitivity}
\label{sec:s-sensitivity}

\begin{table}[h]
\centering
\caption{}
% Threshold rule: selected on training photos vs fixed 0.55, and the whole grid;
% gate exponent gamma grid. Raw mIoU.
% SRC: runs/40, /41, /42 report every threshold; gamma grid in runs/37-41.
\label{tab:s-sensitivity}
\footnotesize\setlength{\tabcolsep}{3pt}
\begin{tabular}{@{}llccc@{}}
\toprule
 & & & \multicolumn{2}{c}{COCO val2017} \\
\cmidrule(lr){4-5}
Factor & Setting & VOC val & VOC cl. & other 60 \\
\midrule
\multirow{2}{*}{Threshold} & selected on train & \x & \x & \x \\
                           & fixed 0.55        & \x & \x & \x \\
\midrule
\multirow{4}{*}{$\gamma$}  & 0.5 & \x & \x & \x \\
                           & 1   & \x & \x & \x \\
                           & 1.5 & \x & \x & \x \\
                           & 2   & \x & \x & \x \\
\midrule
\multirow{2}{*}{Timestep}  & \x  & \x & \x & \x \\
                           & \x  & \x & \x & \x \\
\bottomrule
\end{tabular}
\end{table}

\section{Evaluation subsets}
\label{sec:s-subsets}

\begin{table}[h]
\centering
\caption{}
% Robustness of the main contrasts: primary items (>=1% area) vs all items;
% person class excluded. Differences with 95% intervals.
% SRC: runs/40 and runs/42 secondary outputs.
\label{tab:s-subsets}
\footnotesize\setlength{\tabcolsep}{3pt}
\begin{tabular}{@{}lccc@{}}
\toprule
Subset & VOC val & COCO, VOC cl. & COCO, other 60 \\
\midrule
Primary items ($\geq$1\% area) & \x & \x & \x \\
All items                      & \x & \x & \x \\
Person excluded                & \x & \x & -- \\
\bottomrule
\end{tabular}
\end{table}

\section{Cost}
\label{sec:s-cost}

\begin{table}[h]
\centering
\caption{}
% Runtime and memory per image and class, and one-off cost of building
% reference profiles; ours vs baselines on the same GPU.
% SRC: NEW measurement on one GPU type.
\label{tab:s-cost}
\footnotesize
\begin{tabular}{@{}lccc@{}}
\toprule
Method & s / (image, class) & GPU memory & One-off cost \\
\midrule
\x & \x & \x & \x \\
\method{} & \x & \x & \x \\
\bottomrule
\end{tabular}
\end{table}

\section{Additional qualitative results and failure cases}
\label{sec:s-qualitative}

\begin{figure}[h]
\centering
\figbox{Gallery: more seen and unseen classes; failure cases (ambiguous words, small objects, co-occurring classes).}{5cm}
% FILE: figures/supp_gallery.pdf
\caption{}
\label{fig:s-gallery}
\end{figure}

\end{document}
